% =====================================================================
%  anexo-d-pki.tex — Scripts de la autoridad certificadora interna
% =====================================================================
\chapter{Scripts de la autoridad certificadora interna}\label{anexo:pki}

La autoridad certificadora interna del prototipo y su autoridad de sellado de tiempo local se operan con cinco scripts de bash y dos archivos de configuración de OpenSSL, todos versionados en el directorio \texttt{pki/} del repositorio y probados con OpenSSL 3.5.7. Los scripts comparten tres disciplinas: se ejecutan con \texttt{set -euo pipefail} (cualquier fallo intermedio aborta), verifican que \texttt{openssl} esté disponible antes de operar, y leen el directorio de trabajo de la CA de la variable de entorno \texttt{PKI\_CA\_DIR} (por defecto \texttt{pki-ca} bajo el directorio actual). El directorio de trabajo es un artefacto de tiempo de ejecución: contiene las llaves privadas y la base de datos de la autoridad, se crea con permisos restrictivos (\texttt{umask 077} al generar llaves, \texttt{chmod 700} sobre \texttt{private/}) y jamás se versiona; como red de seguridad, el \texttt{.gitignore} de la raíz ignora \texttt{pki/**/*.key} y \texttt{pki/**/*.pem}. El modelo de confianza es una jerarquía de una sola capa: una raíz autofirmada firma directamente certificados de entidad final, la CRL y el certificado de la TSA, conforme a X.509 y RFC 5280 \cite{cooper2008}.

Estos mismos scripts son los que ejecuta el prototipo: el adaptador de infraestructura de la autoridad los invoca como subprocesos desde los subcomandos \texttt{pki} de la interfaz de línea de comandos (véase el Anexo~\ref{anexo:cli}), y la suite de pruebas los usa como fixtures reales en directorios temporales, de modo que las pruebas ejercitan la misma implementación de autoridad local utilizada por la demostración. Esto no acredita su operación como autoridad de producción o prestador autorizado.

\section{Inventario y parámetros criptográficos comunes}\label{sec:pki-inventario}

\begin{table}[H]
  \centering
  \caption{Archivos del directorio \texttt{pki/} y su propósito.}
  \label{tab:pki-inventario}
  \begin{tabularx}{\linewidth}{@{}l X@{}}
    \toprule
    \textbf{Archivo} & \textbf{Propósito} \\
    \midrule
    \texttt{init-ca.sh} & Crea la CA raíz: llave privada, certificado autofirmado y base de datos de la autoridad. \\
    \texttt{issue-cert.sh} & Emite un certificado de entidad final firmado por la CA para el common name dado. \\
    \texttt{revoke.sh} & Marca como revocado un certificado emitido, en la base de datos de la CA. \\
    \texttt{gen-crl.sh} & Genera la lista de certificados revocados (CRL) en formato PEM. \\
    \texttt{issue-tsa-cert.sh} & Emite el certificado de la TSA local y prepara su directorio de trabajo. \\
    \texttt{openssl.cnf} & Configuración de la CA: layout del directorio, política de emisión y extensiones X.509 v3. \\
    \texttt{tsa.cnf} & Configuración de la TSA para \texttt{openssl ts -reply} y extensiones del certificado de sellado. \\
    \texttt{README.md} & Guía de operación: orden de uso, ejemplo completo y comandos de verificación independiente. \\
    \bottomrule
  \end{tabularx}
\end{table}

Los parámetros criptográficos son idénticos en toda la jerarquía y están alineados con el nivel de fortaleza de 128 bits adoptado por el diseño criptográfico (\cref{sec:nivel-seguridad}): todas las llaves son RSA de 3\,072 bits generadas con \texttt{openssl genpkey} y todas las firmas usan SHA-256, con \texttt{sha256WithRSAEncryption} como algoritmo de firma. La \cref{tab:pki-parametros} resume los valores exactos.

\begin{table}[H]
  \centering
  \caption{Parámetros criptográficos y de vigencia de la autoridad interna.}
  \label{tab:pki-parametros}
  \begin{tabularx}{\linewidth}{@{}X l@{}}
    \toprule
    \textbf{Parámetro} & \textbf{Valor} \\
    \midrule
    Algoritmo y tamaño de llave (raíz, entidad final y TSA) & RSA 3\,072 bits \\
    Digesto de firma en certificados, CRL y sellos & SHA-256 \\
    Vigencia del certificado raíz autofirmado & 1\,825 días (5 años) \\
    Vigencia de los certificados de entidad final & 365 días (1 año) \\
    Vigencia del certificado de la TSA & 365 días (1 año) \\
    Vigencia de cada CRL (\texttt{default\_crl\_days}) & 7 días \\
    Retrodatación del inicio de vigencia (\texttt{notBefore}) & 5 minutos \\
    Serial inicial de certificados y de números de CRL & \texttt{1000} \\
    Serial inicial de tokens de la TSA & \texttt{1000} \\
    Política de sellado de la TSA (OID) & \texttt{1.3.6.1.4.1.13762.3} \\
    Exactitud del tiempo aseverado en cada sello & 1 segundo \\
    \bottomrule
  \end{tabularx}
\end{table}

El sujeto por defecto (sección \texttt{dn\_defaults} de \texttt{openssl.cnf}) es una identidad de prototipo para un despacho jurídico mexicano: \texttt{C=MX}, \texttt{ST=Ciudad de Mexico}, \texttt{L=Ciudad de Mexico}, \texttt{O=Despacho Juridico Demo}, con \texttt{OU} y \texttt{CN} propios de cada pieza (autoridad, personal del despacho o autoridad de sellado).

\section{init-ca.sh: creación de la CA raíz}\label{sec:pki-init-ca}

El script crea el layout completo del directorio de trabajo (\texttt{certs/}, \texttt{crl/}, \texttt{csr/}, \texttt{newcerts/} y \texttt{private/} con permisos 700) e inicializa la base de datos de la autoridad: \texttt{index.txt} vacío, \texttt{serial} en \texttt{1000} y \texttt{crlnumber} en \texttt{1000}. Después genera la llave privada RSA 3\,072 bajo \texttt{umask 077} y autofirma el certificado raíz por 1\,825 días con las extensiones \texttt{v3\_root\_ca}. El script es idempotente y defensivo: sobre una CA ya inicializada no hace nada y lo reporta; rehúsa sobrescribir llave o certificado existentes; y un estado inconsistente ---existe la llave pero no el certificado, o viceversa--- es un error explícito.

\begin{lstlisting}[language=bash, caption={Generación de la llave raíz y autofirma del certificado en \texttt{pki/init-ca.sh}.}, label={lst:pki-init-ca}]
# Generate the CA private key with restrictive permissions.
(
    umask 077
    openssl genpkey -quiet -algorithm RSA \
        -pkeyopt "rsa_keygen_bits:$KEY_BITS" -out "$CA_KEY"
)

# Self-sign the root certificate. The subject comes from the
# dn_defaults section of openssl.cnf.
openssl req -config "$OPENSSL_CNF" -new -x509 -sha256 \
    -days "$CA_DAYS" -extensions v3_root_ca \
    -key "$CA_KEY" -out "$CA_CERT"
\end{lstlisting}

\section{issue-cert.sh: emisión de certificados de entidad final}\label{sec:pki-issue-cert}

Recibe como único argumento el common name del sujeto, restringido por una expresión regular a un subconjunto ASCII seguro (letras, dígitos, espacios y \texttt{. \_ @ -}, iniciando con letra o dígito) para poder incrustarlo en el argumento \texttt{-subj} y en nombres de archivo sin escape. Del common name deriva un identificador en minúsculas con los espacios convertidos a guiones (\texttt{Juan Perez} $\rightarrow$ \texttt{juan-perez}) que nombra la llave, el CSR y el certificado. Genera la llave RSA 3\,072, crea la solicitud de firma y la firma con la CA por 365 días usando las extensiones \texttt{v3\_end\_entity}; el país y la organización del sujeto deben coincidir con los de la CA porque la política de emisión así lo exige. El script rehúsa reemitir si ya existen el certificado o la llave del mismo identificador.

Una decisión operativa merece nota: el inicio de vigencia se retrodata cinco minutos, como acostumbran las autoridades públicas, para que el certificado sea utilizable de inmediato aun cuando la máquina emisora y la máquina que valida difieran ligeramente en su reloj; sin la retrodatación, una validación efectuada en el mismo segundo de la emisión puede caer antes de \texttt{notBefore} y reportar el certificado como aún no válido.

\begin{lstlisting}[language=bash, caption={Retrodatación del inicio de vigencia y firma por la CA en \texttt{pki/issue-cert.sh}.}, label={lst:pki-issue-cert}]
# Backdate the start of validity by a few minutes, as public authorities
# commonly do, so the certificate is immediately usable even when the
# issuing and the relying machine disagree slightly on the time (clock
# skew). Without this, a validation performed in the same second as the
# issuance can land before notBefore and report the certificate as not
# yet valid.
START_DATE="$(date -u -d '5 minutes ago' +%Y%m%d%H%M%SZ)"

# Sign the request with the CA using the end-entity extensions.
openssl ca -config "$OPENSSL_CNF" -batch -notext -md sha256 \
    -startdate "$START_DATE" -days "$CERT_DAYS" -extensions v3_end_entity \
    -in "$CSR_FILE" -out "$CERT_FILE"
\end{lstlisting}

\section{revoke.sh: revocación}\label{sec:pki-revoke}

Recibe la ruta del certificado a revocar, lee su número de serie con \texttt{openssl x509 -noout -serial} y consulta la base de datos \texttt{index.txt} de la CA ---un archivo separado por tabuladores con estado, expiración, fecha de revocación, serial, archivo y sujeto, donde el estado \texttt{R} significa revocado--- antes de intentar la revocación. Si el certificado ya está revocado, lo reporta y termina con éxito, de modo que el script puede re-ejecutarse sin riesgo. La revocación se asienta con \texttt{openssl ca -revoke} y solo se hace visible a terceros cuando se regenera la CRL con \texttt{gen-crl.sh}.

\begin{lstlisting}[language=bash, caption={Consulta del estado y revocación en \texttt{pki/revoke.sh}.}, label={lst:pki-revoke}]
# index.txt is tab separated: status, expiry, revocation date,
# serial, file name, subject. Status R means already revoked.
if awk -F '\t' -v s="$SERIAL" '$1 == "R" && $4 == s { found = 1 } END { exit !found }' \
    "$PKI_CA_DIR/index.txt"; then
    printf 'Certificate with serial %s is already revoked (nothing to do).\n' "$SERIAL"
    exit 0
fi

openssl ca -config "$OPENSSL_CNF" -revoke "$CERT_PATH"
\end{lstlisting}

\section{gen-crl.sh: generación de la CRL}\label{sec:pki-gen-crl}

Produce la lista de certificados revocados en formato PEM a partir de la base de datos de la CA. Cada CRL es válida siete días (\texttt{default\_crl\_days = 7} en \texttt{openssl.cnf}), por lo que debe regenerarse al menos semanalmente y después de cada revocación; el validador del prototipo trata una CRL vencida como error duro, nunca como un pase permisivo. La única extensión de la lista es \texttt{authorityKeyIdentifier} (sección \texttt{crl\_ext}). Al terminar, el script reporta las fechas \texttt{lastUpdate} y \texttt{nextUpdate} de la lista recién emitida y cuenta los certificados revocados que contiene.

\begin{lstlisting}[language=bash, caption={Emisión y reporte de la CRL en \texttt{pki/gen-crl.sh}.}, label={lst:pki-gen-crl}]
openssl ca -config "$OPENSSL_CNF" -gencrl -out "$CRL_FILE"

printf 'CRL generated successfully.\n'
printf '  CRL file: %s\n' "$CRL_FILE"
openssl crl -noout -lastupdate -nextupdate -in "$CRL_FILE" | sed 's/^/  /'
REVOKED_COUNT="$(awk -F '\t' '$1 == "R"' "$PKI_CA_DIR/index.txt" | wc -l)"
printf '  Revoked certificates listed: %s\n' "$REVOKED_COUNT"
\end{lstlisting}

\section{issue-tsa-cert.sh: certificado de la autoridad de sellado}\label{sec:pki-issue-tsa}

Emite el certificado de la TSA local y prepara su directorio de trabajo, tomado de la variable \texttt{TSA\_DIR} (por defecto \texttt{pki-tsa}, junto al directorio de la CA). Genera la llave RSA 3\,072 de la TSA, crea el CSR con el sujeto \texttt{OU=Autoridad de Sellado de Tiempo, CN=TSA Interna Despacho Juridico Demo} y lo firma con la CA por 365 días, con la misma retrodatación de cinco minutos, pero con una diferencia esencial: las extensiones no provienen de \texttt{openssl.cnf} sino de \texttt{tsa.cnf} (\texttt{-extfile tsa.cnf -extensions v3\_tsa}), porque RFC 3161 \cite{adams2001} exige que el \texttt{extendedKeyUsage} de un certificado de sellado sea crítico y contenga \texttt{id-kp-timeStamping}. Finalmente prepara los dos archivos que \texttt{openssl ts -reply} lee en cada emisión de token: el archivo \texttt{serial} de tokens (inicia en \texttt{1000}) y \texttt{tsa-chain.pem}, la cadena que se adjunta a cada respuesta, que es una copia del certificado raíz de la CA.

\begin{lstlisting}[language=bash, caption={Firma con extensiones de \texttt{tsa.cnf} y preparación del directorio TSA en \texttt{pki/issue-tsa-cert.sh}.}, label={lst:pki-issue-tsa}]
# Sign the request with the CA. The extensions come from tsa.cnf, not
# from openssl.cnf, so the certificate carries the timestamping
# extendedKeyUsage instead of the end-entity one.
openssl ca -config "$OPENSSL_CNF" -batch -notext -md sha256 \
    -startdate "$START_DATE" -days "$CERT_DAYS" -extfile "$TSA_CNF" -extensions v3_tsa \
    -in "$CSR_FILE" -out "$CERT_FILE"

# Files "openssl ts -reply" reads: the running serial number of issued
# tokens and the chain appended to every response.
[ -f "$SERIAL_FILE" ] || printf '1000\n' > "$SERIAL_FILE"
cp "$PKI_CA_DIR/ca.crt.pem" "$CHAIN_FILE"
\end{lstlisting}

\section{openssl.cnf: configuración de la CA}\label{sec:pki-openssl-cnf}

El archivo define la sección \texttt{internal\_ca} con el layout del directorio de trabajo (resuelto desde \texttt{\$ENV::PKI\_CA\_DIR}, por lo que el archivo no es utilizable de forma independiente de los scripts) y los valores por defecto para la emisión, que son \texttt{default\_md = sha256}, \texttt{default\_days = 365} y \texttt{default\_crl\_days = 7}. Tres directivas endurecen la emisión: \texttt{copy\_extensions = none} garantiza que solo se apliquen las extensiones definidas en el propio archivo y nunca las que el solicitante incluya en su CSR; la política \texttt{policy\_internal} exige que el país y la organización del certificado emitido coincidan con los de la CA (\texttt{match}) y que el common name esté presente (\texttt{supplied}); y \texttt{unique\_subject = no} permite reemitir para un sujeto cuyo certificado anterior fue revocado. La sección \texttt{req} fija \texttt{default\_bits = 3072}, \texttt{default\_md = sha256} y \texttt{string\_mask = utf8only}. Los tres juegos de extensiones X.509 v3 son los siguientes:

\begin{lstlisting}[language={}, caption={Extensiones X.509 v3 definidas en \texttt{pki/openssl.cnf}.}, label={lst:pki-extensiones}]
[ v3_root_ca ]
# Extensions for the self-signed root CA certificate.
basicConstraints      = critical, CA:true
keyUsage              = critical, keyCertSign, cRLSign
subjectKeyIdentifier  = hash

[ v3_end_entity ]
# Extensions for issued end-entity certificates: people or services
# that sign documents and authenticate against internal systems.
basicConstraints        = critical, CA:false
keyUsage                = critical, digitalSignature, nonRepudiation
extendedKeyUsage        = clientAuth, emailProtection
authorityKeyIdentifier  = keyid:always
subjectKeyIdentifier    = hash

[ crl_ext ]
# Extensions attached to every generated CRL.
authorityKeyIdentifier  = keyid:always
\end{lstlisting}

La raíz declara \texttt{CA:true} con uso de llave restringido a firmar certificados y CRL; las entidades finales declaran \texttt{CA:false} con \texttt{digitalSignature} y \texttt{nonRepudiation} críticos ---los usos que sustentan la firma de documentos--- más autenticación de cliente y protección de correo como usos extendidos.

\section{tsa.cnf: configuración de la autoridad de sellado}\label{sec:pki-tsa-cnf}

La sección \texttt{internal\_tsa} gobierna a \texttt{openssl ts -reply}, la orden con la que la TSA local emite cada token; el directorio se resuelve desde \texttt{\$ENV::TSA\_DIR}. Todos los digestos son SHA-256, tanto los aceptados en las solicitudes como el usado para firmar la respuesta. La política estampada en cada token es el OID \texttt{1.3.6.1.4.1.13762.3}: el arco \texttt{1.3.6.1.4.1} es el árbol de empresas privadas de la IANA y el valor es una política de demostración del prototipo, no una política de producción registrada; las solicitudes que pidan cualquier otra política se rechazan. El tiempo aseverado es exacto a un segundo (\texttt{accuracy = secs:1}), las respuestas no garantizan orden entre tokens (\texttt{ordering = no}) y se mantienen mínimas (\texttt{tsa\_name = no}, \texttt{ess\_cert\_id\_chain = no}).

\begin{lstlisting}[language={}, caption={Parámetros de emisión de tokens en \texttt{pki/tsa.cnf}.}, label={lst:pki-tsa-cnf}]
# All digests are SHA-256: the imprint digests accepted in requests and
# the digest used for the response signature.
digests       = sha256
signer_digest = sha256

# Policy identifier stamped into every token. The arc 1.3.6.1.4.1 is
# the IANA private-enterprise tree; the value below is a demo policy
# for this prototype, not a registered production policy. Requests
# asking for any other policy are refused.
default_policy = 1.3.6.1.4.1.13762.3

# The asserted time is accurate to one second and responses carry no
# ordering guarantee between tokens.
accuracy = secs:1
ordering = no
\end{lstlisting}

El mismo archivo define las extensiones \texttt{v3\_tsa} del certificado de la autoridad de sellado, que difieren de las de entidad final en el punto que RFC 3161 exige: el uso extendido de llave \texttt{timeStamping} marcado como crítico, junto con \texttt{basicConstraints = critical, CA:false} y \texttt{keyUsage = critical, digitalSignature}.

\section{Flujo típico de operación}\label{sec:pki-flujo}

El orden de uso es el siguiente (los pasos de la autoridad certificadora están documentados en \texttt{pki/README.md}; el paso de la TSA lo añade \texttt{issue-tsa-cert.sh}):

\begin{enumerate}
  \item \texttt{init-ca.sh}, una sola vez, para crear la raíz y la base de datos de la autoridad.
  \item \texttt{issue-cert.sh "Nombre Apellido"}, una vez por cada firmante que requiera certificado.
  \item \texttt{issue-tsa-cert.sh}, una sola vez, si se empleará el sellado de tiempo local.
  \item \texttt{revoke.sh ruta/al/certificado.crt.pem}, cuando una llave se compromete o el vínculo con el firmante termina.
  \item \texttt{gen-crl.sh}, después de cada revocación y al menos cada siete días, para que los verificadores dispongan siempre de una lista vigente.
\end{enumerate}

En la operación normal del prototipo estos pasos no se ejecutan a mano: el adaptador de la autoridad en el crate de infraestructura invoca \texttt{init-ca.sh}, \texttt{issue-cert.sh}, \texttt{revoke.sh} y \texttt{gen-crl.sh} como subprocesos desde los subcomandos \texttt{pki} del binario, inyectando \texttt{PKI\_CA\_DIR}; y la TSA local invoca \texttt{openssl ts -query} y \texttt{openssl ts -reply} contra \texttt{tsa.cnf} con \texttt{TSA\_DIR} inyectada, serializando las respuestas concurrentes con un candado consultivo sobre un archivo del directorio TSA, porque OpenSSL avanza el archivo de seriales con una secuencia leer-incrementar-escribir no atómica. Las pruebas de integración conducen estos mismos scripts en directorios temporales y contrastan sus productos con la línea de comandos de OpenSSL (\texttt{crates/infrastructure/tests/\allowbreak openssl\_ca\_adapter.rs}, \texttt{certificate\_validation.rs} y \texttt{local\_tsa\_rfc3161.rs}).

Cualquier producto de la autoridad puede auditarse sin el prototipo, con los comandos documentados en \texttt{pki/README.md}: \texttt{openssl verify -CAfile} para el encadenamiento a la raíz, \texttt{openssl verify -crl\_check} ---concatenando raíz y CRL--- para el estado de revocación, y \texttt{openssl crl -noout -text} para inspeccionar la lista; la guía completa de verificación de un paquete de evidencia se presenta en el Anexo~\ref{anexo:verificacion}.
